\par\clearpage
\SourceLeaf{137}\GaussVTheoryEarlyFolio{133}
\noindent Substituirt man diese Werthe, und die obigen von $\delta^0,\delta',\delta''$ in unserer Gleichung indem man die Distanzen in Secunden ausdrückt, so wird sie
\[
0=17556\omega^0+2774\omega'-20377\omega''
\]
oder
\[
\omega''=0{,}86158\omega^0+0{,}13613\omega'
\]
Aus den beobachteten horizontalen Intensitäten in Göttingen und Mailand folgt hienach die für Paris $\omega''=0{,}51696$, fast genau mit dem beobachteten Werthe $0{,}51804$ übereinstimmend.

Uebrigens sieht man leicht, dass, wenn man sich erlauben will, anstatt der Distanzen $P^0P'$ u. s. w. ihre Sinus zu setzen, die obige Formel unmittelbar durch die geographischen Längen und Breiten der Örter ausgedrückt werden kann.

\Article{11}
Die Linie auf der Erdoberfläche, in deren sämmtlichen Punkten $V$ einerlei bestimmten Werth $=V^0$ hat, scheidet, allgemein zu reden, die Theile jener Fläche, in welchen $V$ einen Werth grösser als $V^0$ hat, von denen, wo er kleiner ist. Die horizontale magnetische Kraft in jedem Punkte dieser Linie ist offenbar senkrecht gegen dieselbe, und zwar nach der Seite zu gerichtet, wo die grössern Werthe von $V$ Statt finden. Ist $\mathrm{d}s$ eine unendlich kleine Linie in dieser Richtung, und $V^0+\mathrm{d}V^0$ der Werth von $V$ an deren anderm Endpunkte, so ist $\frac{\mathrm{d}V^0}{\mathrm{d}s}$ die Intensität der horizontalen magnetischen Kraft an dieser Stelle. So wie nun auch hier die Gesammtheit der Punkte, welchen der Werth $V=V^0+\mathrm{d}V^0$ entspricht, eine zweite der ersten unendlich nahe liegende Linie bildet, also aus der ganzen Erdfläche eine \textit{Zone} aussondert, innerhalb welcher die Werthe von $V$ zwischen $V^0$ und $V^0+\mathrm{d}V^0$ liegen, und wo die horizontale Intensität der Breite der Zone verkehrt proportional ist, so wird, wenn man $V$ durch unendlich kleine aber gleiche Stufen von dem kleinsten auf der Erdoberfläche Statt habenden Werthe bis zum grössten sich ändern lässt, die ganze Erdfläche in eine unendlich grosse Anzahl unendlich schmaler Zonen abgetheilt, gegen deren Scheidungslinien die horizontale magnetische Kraft überall normal, und in ihrer Intensität der Breite der Zonen an den betreffenden Stellen verkehrt proportional ist. Den beiden äussersten Werthen von $V$ entsprechen hiebei zwei von den Zonen eingeschlossene Punkte, in welchen die horizontale Kraft $=0$ wird, und
\par\clearpage
\SourceLeaf{138}\GaussVTheoryEarlyFolio{134}
\noindent wo also die ganze magnetische Kraft nur vertical sein kann: diese Punkte heissen die magnetischen Pole der Erde.

Die Scheidungslinien der Zonen sind nichts anderes, als die Schnitte der im 7. Art. betrachteten Flächen mit der Erdoberfläche, während in den Polen nur Berührung Statt findet.

\Article{12}
Die im vorhergehenden Artikel beschriebene Gestaltung des Liniensystems ist eigentlich nur der einfachste Typus, der mancherlei Ausnahmen erleiden könnte, wenn jede mögliche Vertheilung des Magnetismus in der Erde berücksichtigt werden sollte. Wir werden indess hier diesen Gegenstand nicht erschöpfen, sondern zur Erläuterung nur einige Bemerkungen über die Ausnahmsfälle beifügen. zumal da bei der \textit{wirklichen} magnetischen Beschaffenheit der Erde das Liniensystem auf ihrer Oberfläche allerdings jene Gestaltung hat, wenigstens gewiss keine ins Grosse gehende Ausnahmsfälle darbietet, sondern höchstens vielleicht hie und da einen bloss localen.

Von einigen Physikern ist die Meinung aufgestellt, dass die Erde zwei magnetische Nordpole und zwei Südpole habe: es scheint aber nicht, dass vorher der wesentlichsten Bedingung genügt, und eine \textit{präcise} Begriffsbestimmung gegeben sei, was man unter einem magnetischen Pole verstehen wolle. Wir werden mit dieser Benennung jeden Punkt der Erdoberfläche bezeichnen, wo die horizontale Intensität $=0$ ist: allgemein zu reden ist also daselbst die Inclination $=90^0$; es ist aber auch der singuläre Fall (wenn er vorkäme) mit eingeschlossen, wo die ganze Intensität $=0$ ist. Wollte man diejenigen Stellen magnetische Pole nennen, wo die ganze Intensität einen Maximumwerth hat (d. i. einen grössern, als ringsherum in der nächsten Umgebung): so darf man nicht vergessen, dass dies etwas von jener Begriffsbestimmung ganz verschiedenes ist, dass letztere Punkte mit jenen weder dem Orte noch der Anzahl nach einen nothwendigen Zusammenhang haben, und dass es zur Verwirrung führt, wenn ungleichartige Dinge mit einerlei Namen benannt werden.

Sehen wir von der wirklichen Beschaffenheit der Erde ab, und fassen die Frage allgemein auf, so können allerdings mehr als zwei magnetische Pole existiren; es scheint aber noch nicht bemerkt zu sein, dass sobald z. B. zwei Nordpole vorhanden sind, es nothwendig zwischen ihnen noch einen dritten Punkt
\par\clearpage
\SourceLeaf{139}\GaussVTheoryEarlyFolio{135}
\noindent geben muss, der gleichfalls ein magnetischer Pol, aber eigentlich weder ein Nordpol noch ein Südpol, oder, wenn man lieber will, beides zugleich ist.

Zur Aufklärung dieses Gegenstandes ist nichts dienlicher, als die Betrachtung unsers Liniensystems.

\expandafter\def\csname HistoricalNoteText16\endcsname{Diese Ringlinien sind, selbst als unendlich klein angenommen, nicht nothwendig kreisrund, sondern allgemein zu reden elliptisch, und daher die gegen sie normale Richtung der Magnetnadel nicht mit der Richtung nach $P^*$ zusammenfallend, ausser an vier Stellen jedes Ringes. Man kann daher bedeutende Fehler begehen, wenn man den Durchschnitt von zwei verlängerten Compassrichtungen, aus beträchtlichen Entfernungen, ohne Weiteres für $P^*$ annimmt.}
\expandafter\def\csname HistoricalNoteText17\endcsname{Wir conformiren uns hier dem gewöhnlichen Sprachgebrauche, wonach man den von Capitaine Ross festgelegten Punkt mit jenem Namen belegt, obgleich er eigentlich ein Südpol ist, insofern man die Erde selbst wie einen Magnet betrachtet.}
Wenn die Function $V$ in einem Punkte der Erdoberfläche $P^*$ einen Maximumwerth $V^*$ hat, also ringsum kleinere Werthe, so wird einer Reihe von stufenweise abnehmenden Werthen ein System von Ringlinien entsprechen, deren jede alle vorhergehenden und den Punkt $P^*$ einschliesst, und die Richtung der horizontalen magnetischen Kraft oder des Nordpols der Magnetnadel wird auf jeder dieser Ringlinien \textit{nach Innen} gehen\HistoricalNote{16}: dies ist das charakteristische Merkmal eines magnetischen Nordpols{\renewcommand{\thefootnote}{**)}\HistoricalNote{17}}. Man kann offenbar die Ringe so klein, oder die entsprechenden Werthe der Function $V$ so wenig von $V^*$ verschieden annehmen, dass jeder andere gegebene Punkt noch ausserhalb bleibt.

Wir wollen mit $S$ den Inbegriff aller Punkte auf der Erdoberfläche bezeichnen, in welchen der Werth von $V$ grösser ist als eine gegebene Grösse $W$. Offenbar wird $S$ entweder Einen zusammenhängenden Flächenraum bilden, oder mehrere von einander getrennte, und \textit{in} der Begrenzungslinie oder den Begrenzungslinien, welche dieselbe von den übrigen Theilen, wo $V$ kleiner als $W$ ist, scheiden, wird $V=W$ sein. Lässt man $W$ ab- oder zunehmen, so erweitert oder verengt sich jener Flächenraum.

Nehmen wir nun an, $P^{**}$ sei ein zweiter Punkt von ähnlicher Beschaffenheit wie $P^*$, so dass auch in jenem $V$ einen Maximumwerth $=V^{**}$ habe. Da man, nach dem was vorhin bemerkt ist, der Grösse $W$ einen Werth kleiner als $V^*$ und so wenig davon verschieden beilegen kann, das $P^{**}$ ausserhalb desjenigen Stücks von $S$ fällt, in welchem $P^*$ liegt, so wird, wenn man voraussetzt, dass $V^{**}$ nicht kleiner ist als $V^*$ (was erlaubt ist), mithin auch grösser als $W$,
\par\clearpage
\SourceLeaf{140}\GaussVTheoryEarlyFolio{136}
\noindent nothwendig auch $P^{**}$ einem Stück von $S$ angehören: es liegen folglich $P^*$ und $P^{**}$ zwar beide in $S$, aber in getrennten Stücken von $S$.

Offenbar kann man dagegen auch $W$ so klein annehmen, dass $P^*$ und $P^{**}$ in Einem zusammenhängenden Stücke von $S$ liegen, da, wenn man nur $W$ klein genug nimmt, $S$ die ganze Erdfläche umfassen kann.

Lässt man nun $W$ alle Werthe vom ersten zum zweiten stufenweise durchlaufen, so muss einer darunter $V^{***}$ der letzte sein, für welchen $P^*,P^{**}$ noch in getrennten Stücken von $S$ liegen, welche, sobald $W$ von da noch weiter abnimmt, in Ein Stück zusammenfliessen.

Geschieht dieses Zusammenfliessen in Einem Punkte $P^{***}$. so hat die Begrenzungslinie, in welcher $V=V^{***}$ ist, die Gestalt einer Acht, die in jenem Punkte sich selbst kreuzt, und man überzeugt sich leicht, dass daselbst die horizontale Intensität $=0$ sein muss. In der That geschieht jene Kreuzung entweder unter einem messbaren Winkel, oder nicht. Im erstern Fall müsste die horizontale Kraft, wenn sie nicht $=0$ wäre, gegen zwei verschiedene Tangenten normal sein, was absurd ist; im zweiten Falle, wo die beiden Hälften der Acht in $P^{***}$ einander berühren, oder einerlei Tangente haben würden, könnte die gegen diese Tangente normale Kraft nur gegen das Innere der einen Flächenhälfte der Acht gerichtet sein, was einen Widerspruch enthält, da der Werth von $V$ nach beiden Seiten zu wächst; es ist also $P^{***}$ nach unserer Definition ein wahrer magnetischer Pol, aber ein Pol, welcher in Beziehung auf die zunächstliegenden Punkte innerhalb der beiden Öffnungen der Acht wie ein Südpol, in Beziehung auf die ausserhalb liegenden hingegen wie ein Nordpol betrachtet werden muss. Zur Erläuterung dieser Gestaltung des Liniensystems kann die Fig. 1. dienen.

Geschieht das Zusammenfliessen an zwei verschiedenen Stellen zugleich, so gilt von diesen dasselbe, was eben von Einem Punkte bewiesen ist, und man sieht leicht ein, dass sich dann innerhalb des $P^*$ und $P^{**}$ einschliessenden Raumes ein inselförmiger Raum bilden wird, der bei fortwährender Abnahme von $W$ sich immer mehr verengen, und zuletzt nothwendig in einen wahren Südpol auflösen muss.

Ähnliches gilt, wenn das Zusammenfliessen zugleich in drei oder mehrern einzelnen Punkten Statt findet. Geschieht es aber auf einmal in einer ganzen Linie, so muss auch in allen Punkten derselben die horizontale Kaft verschwinden.

\par\clearpage
\SourceLeaf{141}\GaussVTheoryEarlyFolio{137}
Übrigens ist von selbst klar, dass eben so die Annahme von zwei \textit{Südpolen} zugleich das Dasein eines dritten Polpunkts bedingt, welcher weder Südpol noch Nordpol, oder vielmehr beides zugleich ist.

\Article{13}
Aus dem, was im vorhergehenden Artikel entwickelt ist, übersieht man nun leicht. welche Bewandtniss es mit mehrern denkbaren Ausnahmen von dem einfachsten Typus unsers Liniensystems habe. Der Inbegriff aller Punkte, denen ein bestimmter Werth von $V$ entspricht, kann eine Linie sein, die aus mehrern Stücken besteht, wovon jedes in sich selbst zurückkehrt, die aber ganz von einander getrennt sind; es kann eine Linie sein, die sich selbst kreuzt; endlich kann es auch eine solche sein, der auf beiden Seiten Flächenräume anliegen, wo $V$ grösser ist als in der Linie, oder auf beiden Seiten kleiner.

Wir können behaupten, dass etwas ins Grosse gehende Abweichungen solcher Art vom einfachsten Typus auf der Erde nicht Statt finden. Aber locale Abweichungen sind sehr wohl denkbar, wo nahe unter der Erdoberfläche magnetische Massen sich befinden, die zwar in etwas beträchtlicher Entfernung keine merkliche Wirkung mehr ausüben, aber in der unmittelbaren Umgebung doch eine so starke, dass die in regelmässiger Fortschreitung wirkende erdmagnetische Kraft davon ganz überboten und unkenntlich gemacht wird. In der einfachsten Form könnte dann das Liniensystem in einer solchen Gegend eine Gestaltung haben, wie die $2^{\mathrm{te}}$ Figur versinnlicht.

\Article{14}
Nach dieser geometrischen Darstellung der Verhältnisse der horizontalen magnetischen Kraft schreiten wir zur Entwickelung der Art, wie sie dem Calcül unterworfen werden, fort. Auf der Oberfläche der Erde geht $V$ in eine blosse Function zweier veränderlichen Grössen über, wofür wir die geographische Länge von einem beliebigen ersten Meridian östlich gezählt und die Distanz vom Nordpol annehmen wollen: jene soll mit $\lambda$, diese, das Complement der geographischen Breite, mit $u$ bezeichnet werden. Betrachten wir die Erde als aus der Umdrehung einer Ellipse, deren halbe grosse Axe $=R$, die halbe kleine $=(1-\varepsilon)R$, um letztere entstanden, so ist die Grösse eines Elements des Meridians
\par\GaussVTheoryGathering{18}
